\begingroup
\renewcommand{\arraystretch}{1.45}
\clearpage\fancyhead[L]{\small ARITHMETISCH-GEOMETRISCHES MITTEL}
\pdfbookmark[0]{ARITHMETISCH-GEOMETRISCHES MITTEL}{work-native370}
\SourceLeaf{370}
\selectlanguage{ngerman}
\begin{center}
\textbf{NACHLASS.}\par
\textbf{[ARITHMETISCH-GEOMETRISCHES MITTEL.]}\par
\textbf{TEIL I.}\par
\textbf{ÜBER DEN URSPRUNG UND DIE}\\\textbf{ALLGEMEINEN EIGENSCHAFTEN}\par
\textbf{ARITHMETISCH-GEOMETRISCHER MITTEL.}\par
\rule{0.09\linewidth}{0.4pt}\par
\end{center}
\AGMArticle{1}
\Needspace{22\baselineskip}
Es seien $\left\{\begin{array}{l}a,\allowbreak \ a',\allowbreak \ a'',\allowbreak \ a'''\ .\,\allowbreak .\,\allowbreak .\,\allowbreak .\\ b,\allowbreak \ b',\allowbreak \ b'',\allowbreak \ b'''\ .\,\allowbreak .\,\allowbreak .\,\allowbreak .\end{array}\right\}$ zwei Folgen von Größen, die nach der Regel gebildet werden, dass
ihre entsprechenden Glieder \textit{Mittel} zwischen den vorhergehenden Gliedern sind,
und zwar die Glieder der oberen Folge arithmetische, die der unteren geometrische Mittel; also
\[
\begin{aligned}&a'=\tfrac12(a+b),\quad b'=\surd ab,\\[6pt]&a''=\tfrac12(a'+b'),\quad b''=\surd a'b',\\[6pt]&a'''=\tfrac12(a''+b''),\quad b'''=\surd a''b''\ \mathrm{usw.}\end{aligned}
\]
Wir nehmen dabei an, dass $a$, $b$ reell und positiv sind und für die Quadratwurzeln
überall die positiven Werte genommen werden. So können die Folgen beliebig weit fortgesetzt
werden; alle ihre Glieder sind vollständig bestimmt und haben positive reelle Werte. Zunächst ergeben sich unmittelbar die folgenden Beobachtungen:
\par
I. Ist $a=b$, so sind alle Glieder beider Folgen $=a=b$.
\par
II. Sind dagegen $a$, $b$ ungleich, so ist $(a'-b')(a'+b')=\tfrac14(a-b)^2$; daraus
folgt $b'<a'$, und ebenso ist $b''<a''$, $b'''<a'''$ usw., das heißt, jedes Glied der unteren Folge ist kleiner als das entsprechende Glied der oberen. Daher
nehmen wir in diesem Fall auch $b<a$ an.
\par
\Needspace{13\baselineskip}
III. Unter derselben Voraussetzung gilt $a'<a$, $b'>b$; $a''<a'$, $b''>b'$ usw.; die obere
Folge nimmt also beständig ab, die untere beständig zu. Daraus ist ersichtlich,
dass beide einen Grenzwert haben; diese Grenzwerte lassen sich zweckmäßig mit $a^{\infty}$, $b^{\infty}$ bezeichnen.
\par
\Needspace{19\baselineskip}
IV. Schließlich folgt aus $\frac{a'-b'}{a-b}=\frac{(a-b)}{4(a'+b')}=\frac{a-b}{2(a+b)+4b'}$ die Ungleichung $a'-\nobreak{}b'<\nobreak{}\tfrac12(a-\nobreak{}b)$;
ebenso ist $a''-b''<\tfrac12(a'-b')$ usw. Daraus folgt, dass $a-b$, $a'-b'$,
\SourceLeaf{371}
$a''-b''$, $a'''-b'''$ usw. eine beständig abnehmende Folge bilden, deren
Grenzwert $=0$ ist. Somit gilt $a^{\infty}=b^{\infty}$, das heißt, die obere und die untere Folge haben denselben
Grenzwert, wobei die erste stets größer, die zweite stets kleiner bleibt als dieser.
\par
Diesen Grenzwert nennen wir das \textit{arithmetisch-geometrische Mittel zwischen $a$ und $b$}
und bezeichnen ihn mit $\mathrm{M}(a,b)$.
\Needspace{35\baselineskip}
\AGMArticle{2}
Die Wurzeln der Gleichung $xx-2ax+bb=0$ sind reell und positiv, sofern
$a\geqq b$ gilt; ihr arithmetisches Mittel ist $a$, ihr geometrisches $b$. Bezeichnet
man also eine Wurzel, und zwar die größere, wenn sie ungleich sind, mit ${}^{\prime}\!a$, die andere mit ${}^{\prime}\!b$, so lässt sich ${}^{\prime}\!a$ als das Glied der oberen Folge vor $a$ auffassen, und ebenso ${}^{\prime}\!b$ als das Glied der unteren Folge vor $b$. Bezeichnet man in gleicher Weise die Wurzeln wie folgt,
\begingroup\fontsize{9}{11}\selectfont
\begingroup\setlength{\arraycolsep}{2pt}
\[
\begin{array}{@{}lllll@{}}
\mathrm{Gleichung}&xx-2\,{}^{\prime}\!ax+{}^{\prime}\!b\,{}^{\prime}\!b=0&\mbox{größere Wurzel mit}&{}^{\prime\prime}\!a&\mathrm{kleinere\ mit}\ {}^{\prime\prime}\!b\\[6pt]
&xx-2\,{}^{\prime\prime}\!ax+{}^{\prime\prime}\!b\,{}^{\prime\prime}\!b=0&&{}^{\prime\prime\prime}\!a&\hfill{}^{\prime\prime\prime}\!b\\[6pt]
&xx-2\,{}^{\prime\prime\prime}\!ax+{}^{\prime\prime\prime}\!b\,{}^{\prime\prime\prime}\!b=0&&{}^{\prime\prime\prime\prime}\!a&\hfill{}^{\prime\prime\prime\prime}\!b
\end{array}
\]\endgroup
\endgroup
so lassen sich ${}^{\prime}\!a$, ${}^{\prime\prime}\!a$, ${}^{\prime\prime\prime}\!a$ usw. als Fortsetzung der oberen Folge nach links auffassen und ${}^{\prime}\!b$, ${}^{\prime\prime}\!b$, ${}^{\prime\prime\prime}\!b$ usw. als Fortsetzung der unteren. Damit
liegen nun zwei Folgen vor, die sich in beiden Richtungen ins Unendliche fortsetzen lassen:
\[
\begin{array}{@{}ll@{}}
.\,.\,.\,.\ {}^{\prime\prime\prime\prime}\!a,\;{}^{\prime\prime\prime}\!a,\;{}^{\prime\prime}\!a,\;{}^{\prime}\!a,\; a,\; a',\; a'',\; a''',\; a''''\ .\,.\,.\,.\,.&\mathrm{(I)}\\[6pt]
.\,.\,.\,.\ {}^{\prime\prime\prime\prime}\!b,\;{}^{\prime\prime\prime}\!b,\;{}^{\prime\prime}\!b,\;{}^{\prime}\!b,\; b,\; b',\; b'',\; b''',\; b''''\ .\,.\,.\,.\,.&\mathrm{(II)}
\end{array}
\]
\Needspace{32\baselineskip}
\par Jedes Glied der Folge (I) ist also größer als das entsprechende Glied der Folge (II); die erste nimmt von links nach rechts beständig ab und von rechts nach links
zu; die zweite nimmt von links nach rechts beständig zu und in der entgegengesetzten Richtung ab.
Nach rechts haben beide Folgen denselben Grenzwert; nach links dagegen wächst (I)
über jede Schranke hinaus, während (II) den Grenzwert $0$ hat, sofern nicht alle Glieder beider Folgen gleich sind. Denn ${}^{\prime}\!a=a+\surd(aa-bb)$; ${}^{\prime}\!b=a-\surd(aa-bb)$;
daraus folgt ${}^{\prime}\!a\,{}^{\prime}\!a-{}^{\prime}\!b\,{}^{\prime}\!b=4a\surd(aa-bb)>4(aa-bb)$, und ebenso ${}^{\prime\prime}\!a\,\allowbreak {}^{\prime\prime}\!a-{}^{\prime\prime}\!b\,\allowbreak {}^{\prime\prime}\!b>4({}^{\prime}\!a\,\allowbreak {}^{\prime}\!a-{}^{\prime}\!b\,\allowbreak {}^{\prime}\!b)$
usw. Somit ist ersichtlich, dass die Folge $aa-bb$, ${}^{\prime}\!a\,{}^{\prime}\!a-{}^{\prime}\!b\,{}^{\prime}\!b$, ${}^{\prime\prime}\!a\,{}^{\prime\prime}\!a-{}^{\prime\prime}\!b\,{}^{\prime\prime}\!b$ usw. und folglich auch die Folge
$a$, ${}^{\prime}\!a$, ${}^{\prime\prime}\!a$ usw. jede Schranke überschreiten kann; ferner gilt $\frac{{}^{\prime}\!b}{b}=\frac{b}{{}^{\prime}\!a}<\frac{b}{a}$ und ebenso $\frac{{}^{\prime\prime}\!b}{{}^{\prime}\!b}<\frac{{}^{\prime}\!b}{{}^{\prime}\!a}$ usw.,
das heißt, $\frac{{}^{n+1}\!b}{{}^{n}\!b}$ lässt sich durch Vergrößerung von $n$ unter jede Schranke bringen; daher ist der Grenzwert
der Folge $b$, ${}^{\prime}\!b$, ${}^{\prime\prime}\!b$ usw.\nobreak{} $=0$.
\Needspace{9\baselineskip}
\SourceLeaf{372}
\par 
Aus der Definition des arithmetisch-geometrischen Mittels ist der folgende Satz so ersichtlich,
dass es nun keiner weiteren Erläuterung bedarf:

\textsc{Satz.}\quad \textit{Das Mittel zwischen beliebigen entsprechenden Gliedern der Folgen I, II ist dasselbe wie zwischen $a$ und $b$.}

\clearpage
\endgroup
